\section{Executable residual certificates}
\label{sec:executable-certificates}

An asymptotic remainder describes a limiting scale; it need not supply a
constant or a numerical neighborhood. A root obtained by numerical iteration
is useful independent evidence, but is not a verified enclosure.
The function \wl{InverseCertificate} implements the residual bracket of
Theorem~\ref{thm:bracket} on an explicit source interval. Its successful output
contains rational endpoints enclosing a unique root, a bound for the error of
the reported center, and the derivative and residual enclosures used to prove
those statements.

\subsection{The certificate theorem with interval data}

Let $I=[A,B]$, let $c\in(A,B)$, and suppose a structural interval evaluator
proves that $F$ is continuous on $I$ and differentiable inside, with
\[
 F'(I)\subseteq[d_-,d_+],\qquad
 F(c)-y\in[\rho_-,\rho_+].
\]
Assume either $d_->0$ or $d_+<0$, and put
\[
 \mu=\min(|d_-|,|d_+|),\quad
 M=\max(|d_-|,|d_+|),\quad
 \epsilon=\max(|\rho_-|,|\rho_+|),\quad h=\epsilon/\mu.
\]
If $[c-h,c+h]\subseteq I$, Theorem~\ref{thm:bracket}, after changing the sign
of $F$ if necessary, proves a unique root $x_*$ in $I$ and
\[
 |x_*-c|\le h.
\]
The mean-value identity further gives
\[
 x_*\in I\cap[c-h,c+h]\cap
 \left(c-\frac{[\rho_-,\rho_+]}{[d_-,d_+]}\right).
\]
Division is interval division; its denominator is separated from zero. This
last interval is the returned \wl{"RootEnclosure"}. If the residual interval
does not contain zero, there is also the lower bound
\[
 |x_*-c|\ge\frac{\min(|\rho_-|,|\rho_+|)}{M}.
\]
It can prove that a requested accuracy is impossible for a fixed center.

If the symmetric residual bracket does not fit inside $I$, the evaluator
also encloses the residual at $A$ and $B$. Proved opposite signs in the
orientation of $F'$ establish existence by continuity, and the separated
derivative still establishes uniqueness. Once existence in $I$ is known,
the mean-value theorem gives the same bound $|x_*-c|\le h$ and the same
intersection enclosure, without requiring the symmetric bracket to be
contained in $I$. The output identifies this alternative as exact endpoint
sign evidence. This case matters for a coarse asymptotic seed close to one
end of a valid interval: increasing arithmetic precision cannot repair
an unsuitable symmetric-bracket geometry.

The theorem asserts uniqueness on the supplied interval, which the evaluator
also checks lies on the requested side of the source endpoint. It does not
certify global injectivity or identify disconnected monotone intervals as the
same continuation of an inverse germ. This scope is recorded in the result.

\subsection{Rational arithmetic and elementary enclosures}

The proof path uses exact rational endpoints. Addition, multiplication,
reciprocals, and integer powers use the usual enclosing interval operations.
After arithmetic, an endpoint $q\ne0$ is rounded outward to a dyadic grid with
a prescribed number of significant bits. Its magnitude exponent is determined
from the binary lengths of the exact numerator and denominator, followed by
one rational comparison. Rounding uses $\lfloor q/\eta\rfloor\eta$ at a lower
endpoint and $\lceil q/\eta\rceil\eta$ at an upper endpoint. Thus even rounding
is justified by integer arithmetic. Approximate decimal values and
floating-point comparisons are absent from every enclosure and every final
sign or containment test.

For $0\le t\le\tfrac12$, the exponential evaluator uses
\[
 S_N(t)=\sum_{j=0}^N\frac{t^j}{j!},\qquad
 0\le e^t-S_N(t)\le
 \frac{t^{N+1}}{(N+1)!}\frac1{1-t/(N+2)}.
\]
The bound follows by majorizing the ratio of successive omitted terms by
$t/(N+2)$. A positive argument is reduced by powers of two, and the enclosure
is squared back using directed rounding. Negative arguments use the reciprocal
of the positive exponential enclosure. An explicit resource bound on the
argument prevents unbounded integer growth.

For $1\le q\le2$, put $v=(q-1)/(q+1)\in[0,\tfrac13]$. Then
\[
 T_N(q)=2\sum_{j=0}^{N-1}\frac{v^{2j+1}}{2j+1},\qquad
 0\le\log q-T_N(q)\le
 \frac{2v^{2N+1}}{(2N+1)(1-v^2)}.
\]
Every positive rational argument is written exactly as $2^k q$ with
$1\le q<2$, using binary numerator and denominator lengths. The identity
$\log(2^kq)=k\log2+\log q$ therefore reduces all logarithms to the displayed
bound. Both tails are positive rational expressions. Increasing $N$ improves
their analytic precision independently of the directed arithmetic precision.

Before enclosing a logarithm, exact structural identities can remove an
exponential whose real exponent has been enclosed, or split a product whose
factors each have a real logarithm enclosure. Positivity is therefore checked
before distributing logarithms. If, for example, a positive product has
negative factors, the evaluator retains the direct interval route. This avoids
materializing a huge exponential only to take its logarithm while preserving
the real-branch conditions.

Noninteger powers on a strictly positive base are evaluated through
$a^b=\exp(b\log a)$, including supported exact irrational exponents. A positive
rational power also admits a zero lower base endpoint. Expressions are handled
by structural recursion through sums, products, powers, exponentials and
logarithms. A reciprocal crossing zero, a logarithm reaching a nonpositive
argument, or an unsupported head produces a failure. In particular, a
numerical value of an unsupported expression is never accepted as a proof.
Successful structural evaluation establishes the needed continuity on the
whole interval; successful evaluation of the symbolic derivative establishes
the derivative enclosure. Symbolic parameters without exact supported values
are currently rejected with the unresolved expression recorded.

\subsection{Interface, phase coordinates, and adaptive accuracy}

The basic call is
\begin{lstlisting}
s = AsymptoticInverse[x + x^2, {x, 0}, {y, 3}];
InverseCertificate[s, 6, "Interval" -> {1, 3},
  "Center" -> 21/10]
\end{lstlisting}
The independently known root is $2$. The derivative enclosure is $[3,7]$;
the certified bracket encloses $2$ even though the asymptotic expansion was
not requested at a target where its truncation is especially accurate. The
residual theorem is independent of the formal truncation theorem.

Targets must be exact supported numeric expressions. Interval endpoints and
an explicitly supplied center must be rational, with the center strictly
inside the interval. With \wl{"Center" -> Automatic}, numerical evaluation of
the expansion chooses a rational center; an unusable seed is replaced by the
interval midpoint. This decimal calculation selects a candidate only. All
proofs are subsequently rebuilt with exact rational arithmetic. The returned
error bound applies to that recorded center, including its rationalization.
It does not silently claim an error bound for a different displayed expression.

The equation being certified is always the object's stored explicit
\wl{"Function"}. A declared \wl{"InputRemainder"} describes further possible
forward functions, but supplies no computable interval bound for their values
or derivatives. The certificate therefore does not cover those unspecified
terms. Results retain that descriptor, set
\wl{"CertifiesInputRemainderFamily" -> False}, and label the scope
\wl{"StoredExpressionOnly"} when it is nontrivial. An automatically computed
forward truncation can also populate the remainder field, so this label does
not guess the descriptor's provenance. The original function and exact target
are recorded separately from any phase equation. Asymptotic remainder data
are never silently promoted into a numerical error certificate.

For an object using a logarithmic target transformation, or a recognized
Lambert exponential core admitting the same exact transformation, the certificate
evaluates the exact phase and transformed target. The source-domain condition
of the coordinate map is proved on the entire interval first. Since the
exponential is strictly monotone there, the resulting root is a root of the
original equation. This avoids constructing enormous values such as
$\exp(x^2+x)$ merely to subtract two nearly equal numbers.

The option \wl{"TargetError" -> epsilon} requests an exact positive rational
tolerance. After an initial certificate, a sharper interval-Newton enclosure
can supply a new center and verification interval. The preceding certificate
already proves a root exists in that interval, so subsequent error estimates
need not re-prove existence by the initial bracket-containment condition.
Derivative bounds and residual bounds are nevertheless recomputed. Unless
disabled by \wl{"RefineExpansion" -> False}, refinement also tries a larger
asymptotic term goal as a source of a better seed; every accepted seed is
verified independently. A bounded failed or unsupported expansion refinement
does not invalidate an already established interval theorem.

An \wl{AsymptoticCoreInverse} object is also accepted as a seed. Its exact
Lambert core and marker corrections may be evaluated numerically to select
the center, while the certificate evaluates the full explicit
\wl{"Core" + "Perturbation"} equation with rational enclosures. If expansion
refinement is requested, the stored exact core inverse is replayed at a larger
marker depth. Its symbolic identity check and computation are time bounded;
a failed replay falls back to the already certified interval information.
The marker majorant's uncomputed constants are not used as numerical bounds.

With an explicit fixed center, refinement increases enclosure precision without
silently changing the center. A certified lower error bound above the target
causes \wl{Failure["AccuracyFloor", ...]}. Exhausting the refinement budget
causes \wl{Failure["AccuracyNotReached", ...]}, retaining the best established
certificate if one exists. Neither outcome is reported as a successful accuracy
request. The \wl{"History"} property records the attempted centers and
enclosure orders. Resource failures and unresolved derivative signs are
reported separately from mathematical impossibility.

\paragraph{Relative accuracy.}

An optional \wl{"RelativeError" -> tau} requests a relative tolerance with
an independently proved root magnitude. If the root interval $I$ is
separated from zero, set $m_I=\min_{z\in I}|z|$ and
$M_I=\max_{z\in I}|z|$. An absolute error bound $r$ proves the requested
combined tolerance when
\[
 r\le\max(\epsilon,\tau m_I),
\]
where $\epsilon$ is an explicitly supplied \wl{TargetError}, or zero when
none is supplied. This implies
$|A-x_*|\le\max(\epsilon,\tau|x_*|)$ for the unique enclosed root.
The reported relative error bound is $r/m_I$. A proved zero root requires
an explicit positive absolute fallback; no relative quotient at zero is
invented. If an interval merely crosses zero, further refinement can
establish a nonzero lower magnitude or report that the goal was not reached.
For a fixed-center impossibility claim, the opposite bound is needed:
the certified lower error must exceed $\max(\epsilon,\tau M_I)$.
Using $m_I$ in that last test would incorrectly turn a conservative
success criterion into an impossibility assertion.

\subsection{Validation scope}

The regression cases exercise outward rounding at both signs and widely
different magnitudes, explicit exponential and logarithmic bounds, interval
powers crossing zero, positive and negative derivatives, and a zero residual.
Known rational roots provide independent checks of the final enclosures.
Further cases cover the lower Lambert branch, a logarithmic phase, successful
adaptive accuracy, a proved fixed-center accuracy floor, and rejection of
rounded targets, turning points, poles, invalid logarithm domains, and
unsupported heads. These tests validate an executable exact-arithmetic
implementation of the theorem; they do not constitute a separate formal
verification of the Wolfram evaluator.

The package's two motivating functions
$x+x^2(1+\log x)$ and $x+x^{\sqrt2}$ are included, together with a nonunit
leading power and coefficient, both real Lambert branches, and exact-core
perturbations. A supplied interval immediately below $x=e^{-1}$ tests the
small derivative near the turning point of $x\log x$. There the reported
derivative lower bound is small but strictly positive in absolute value; the
result remains an interval theorem, without relying on a logarithmic
asymptotic truncation being accurate near that threshold.
